\documentclass[10pt,a4paper]{amsart}
\usepackage[margin=19mm]{geometry}
\usepackage{amsmath,amssymb,mathtools}
\usepackage[scr=rsfs,cal=euler]{mathalfa}
\usepackage{xfrac}
\usepackage[unicode,hidelinks,bookmarks=false]{hyperref}
\newcommand{\ie}{i.e.\ }
\newcommand{\cf}{cf.\ }
\newcommand{\resp}{respectively\ }
\newcommand{\Subsection}{\subsection}
\providecommand{\nobreakdash}{\nobreak\hbox{-}}
\setlength{\emergencystretch}{2em}
\multlinegap0pt
\usepackage{amssymb}
\newtheorem{lemma}{Lemma}
\title{A determinant on birational maps of Severi-Brauer surfaces}
\author{Elias Kurz}
\date{Source: arXiv:2502.02981v1 (CC BY 4.0)}
\begin{document}
\maketitle

\noindent\emph{Source and licence.} A determinant on birational maps of Severi-Brauer surfaces. Version arXiv:2502.02981v1 (CC BY 4.0), distributed by arXiv under the Creative Commons Attribution 4.0 International licence, http://creativecommons.org/licenses/by/4.0/, as recorded in the arXiv licence metadata for this exact version and shown on the abstract page https://arxiv.org/abs/2502.02981v1. Full licence notice: \texttt{licenses/arxiv-2502.02981-CC-BY-4.0.txt}.\par\smallskip

\noindent\textit{Editorial scope.} Counted unit: the article's lemma LargeCalc (source0.tex lines 423--428). The article's complete original proof of it is delivered with it (source0.tex lines 429--465, one \texttt{proof} environment of 37 lines). No mathematics is written, repaired or re-derived: the statement and the proof are the article's own lines, in the article's own notation, and no retained line is substituted or re-flowed.  No further source-local context is needed: every cross-reference of every retained line resolves inside the two delivered spans.  Nothing was omitted from any retained span, so the package carries no omission record.  No retained line contains a citation, so the wrapper carries no bibliography and no bibliography entry.  No retained line contains a figure environment, an image-inclusion command or drawing code, so no image is delivered and the package declares no asset dependency.  Editorial formatting only, in addition: an AMS \texttt{amsart} preamble that restates the article's own declarations and macros used by the retained lines, namely the environments \texttt{lemma} on separate counters, together with the standard packages those lines need.  The retained environments print in this document as Lemma 1 in order of appearance, whereas the article prints the same items with the numbering of its own full document.  The complete native source of the article as distributed in the arXiv e-print for this exact version is bundled unchanged as \texttt{sources/172775/source0.tex}.
\begin{lemma} \label{LargeCalc}
    Let $\textbf{C}$ be a field and $A \in \mathrm{Gl}_3(C)$ ( where the classes of its columns in $\mathbb{P}^2_C$ together with the coordinate points are in general positions). Then if we define $M_1(A):= A, \, M_i(A) := (M_{i-1}(A)^{\Sigma})^{-1}, \, i=2, \dots , 6$, where $\Sigma = (yz,xz,xy)$, we get, that \[d(A):= det(A^{-1}M_3(A)^{-1}M_5(A)^{-1})det(M_2(A)M_4(A)M_6(A))\] \[= \frac{P(A)^{14}}{P(A^{-1})^7 det((A^{\sigma}))^{42}det(A)^{42}}\] where $P$ is the product of all matrix entries. And as birational maps over $\mathbb{P}^2$ we get, that $M_6(A) \sigma \dots M_2(A) \sigma A \sigma = D(A)$ where $D(A)$ is a diagonal matrix such that:
    \newline
    \[D(A)_{ii} = det(A)^6\frac{\prod_{k=1}^3(A^{-1})_{ik}^2}{\prod_{l=1}^3a_{li}}\]
    Thus $det(D(A)) =\frac{P(A^{-1})^2det(A)^{18}}{P(A)}$.
\end{lemma}

\begin{proof}
    To calculate the determinant we need to calculate $det(M_i(A)), i=2, \dots , 6$. We calculate directly:
    \[det(M_3(A))=-(P(A)det(A)det(M_2(A))^4)^{-1}\]
    With this we can calculate $det(M_i(A)), i \geq 3$ by using that $M_i(A) = M_{i-1}(M_2(A))$ and replacing $det(M_{i-1}(\cdot))$ with the previous equality
    \[det(M_4(A))=-P(M_2(A))^{-1}P(A)^4det(A)^4det(M_2(A))^{15}\]
    \[det(M_5(A))=P(M_3(A))^{-1}P(M_2(A))^4P(A)^{-15}det(A)^{-15}det(M_2(A))^{-56}\]
    \[det(M_6(A))=P(M_4(A))^{-1}P(M_3(A))^4P(M_2(A))^{-15}P(A)^{56}det(A)^{56}det(M_2(A))^{209}\]
    Thus we need to find $P(M_i(A)), i=2,3,4$:
    \[P(M_2(A)) = -P(A^{-1})P(A)^2det(A)^9det(M_2(A))^9\]
    Using that $P(M_2(A)^{-1}) = P(A)^2$ we find using again that $M_i(A) = M_{i-1}(M_2(A))$:
    \[P(M_3(A)) = P(A)^{-3}P(A^{-1})^2det(A)^9det(M_2(A))^{-9}\]
    \[P(M_4(A))=P(A^{-1})^{-3}P(A)^7det(A)^{-18}det(M_2(A))^{18}\]
    Putting this together we get:
    \[det(A^{-1}M_3(A)^{-1}M_5(A)^{-1})det(M_2(A)M_4(A)M_6(A)) = \frac{P(A)^{14}}{P(A^{-1})^7 det((A^{\sigma}))^{42}det(A)^{42}}\]
    The fact that $D(A)$ is an diagonal automorphism, is due to our choice of $M_2(A), \dots , M_6(A)$, which makes sure that the degree is $1$ and $[1:0:0], [0:1:0],[0:0:1]$ are fixed. This follows from the fact that $M_6(A) \sigma \dots \sigma M_1(A)$ is of characteristic $2;1^3$ and it and its inverse have $\{ [1:0:0], [0:1:0],[0:0:1]\}$ as their base points. We then calculate, that that the coordinate lines are contracted onto the base points not on them, which implies that $D(A)$ is diagonal. We also know with a similar argument as above just using the inverse relation, that this diagonal must map the classes of the columns of $(A^{-1})^{\sigma}$ to the classes of the columns of $M_6(A)$. Thus we now calculate the $M_i(A)$ and confirm the announced matrix $D(A)$. We observe, that for $\lambda_i(A) \in \textbf{k}^*, \, i = 2, \dots , 6$:
    \[M_2(A)_{ij} = (\prod_{k \neq i} a_{jk}) (A^{-1})_{ij} \lambda_2(A)\]
    \[\lambda_2(A) = \frac{det(A)}{det(A)^{\Sigma}}\]
    Using $M_i(A) = M_{i-1}(M_2(A))$ and dividing by $P(A), P(A^{-1})$ if needed, we get:
    \[M_3(A)_{ij} = \frac{\prod_{k \neq i}(A^{-1})_{jk}}{\prod_{k = 1}^3a_{ik}} \lambda_3(A)\]
    \[\lambda_3(A) = \lambda_2(M_2(A)) \lambda_2(A)^2P(A)\]
    \[M_4(A)_{ij} = \frac{1}{(\prod_{k \neq i}a_{jk})(\prod_{k=1}^3 (A^{-1})_{ik})} \lambda_4(A)\]
    \[\lambda_4(A) = \frac{\lambda_3(M_2(A))}{\lambda_2(A)}\]
    \[M_5(A)_{ij} = \frac{\prod_{k=1}^3a_{ik}}{\prod_{k \neq i}(A^{-1})_{jk} \prod_{k \neq i} \prod_{l \neq j} a_{kl}} \lambda_5(A)\]
    \[\lambda_5(A) = \frac{\lambda_4(M_2(A))}{P(A) \lambda_2(A)^2}\]
    \[M_6(A)_{ij} = \frac{(\prod_{k=1}^3(A^{-1})_{ik}) (\prod_{k \neq i}a_{jk}^2)a_{ji}}{(\prod_{l \neq j}a_{li}) (\prod_{k \neq i} \prod_{l \neq j}(A^{-1})_{kl})} \lambda_6(A)\]
    \[\lambda_6(A) = \frac{\lambda_5(M_2(A)) \lambda_2(A)}{P(A)}\]
    Now we also find:
    \[D(A)_{ii} = det(A)^6\frac{\prod_{k=1}^3(A^{-1})_{ik}^2}{\prod_{l=1}^3a_{li}}\]
    \[(A^{-1})^{\sigma}_{ij} = \prod_{l \neq i} (A^{-1})_{lj}\]
    Thus:
    \[(D(A)(A^{-1})^{\sigma})_{ij} = det(A)^6\frac{\prod_{k=1}^3(A^{-1})_{ik}^2\prod_{l \neq i} (A^{-1})_{lj}}{\prod_{l=1}^3a_{li}}\]
    and
    \[M_6(A){ij} = \frac{(\prod_{k=1}^3(A^{-1})^2_{ik}) (\prod_{k = 1}^3a_{jk}^2)(\prod_{l \neq i}(A^{-1})_{lj})}{(\prod_{l = 1}^3a_{li})} \frac{\lambda_6(A)}{P(A^{-1})}\]
    Thus for the diagonal matrix $D'(A)$ with $D'(A)_{jj} = (\prod_{j=1}^3a_{jk}^2) \frac{\lambda_6(A)}{P(A)det(A)^6}$ we get:
    \[M_6(A) = D(A) (A^{-1})^{\sigma}D'(A)\]
    Thus the diagonal sends the projective class of the columns of $(A^{-1})^{\sigma}$ to the ones of the columns of $M_6(A)$.
\end{proof}

\end{document}
